\section{Discussion and Future Directions}
\label{sec:discussion}

Synthesizing findings across RQ1 through RQ4 exposes a central structural tension in modern digital forensics. Generative language models accelerate initial triage dramatically, yet their probabilistic execution clashes directly with evidentiary determinism, statutory privacy, and adversarial resilience.

\subsection{The Velocity-Verifiability Paradox}
Empirical studies across triage, incident response, malware classification, and memory analysis confirm tangible operational speedups~\cite{oh2024_volgpt, liu2025_autobnb, clark2025_large}. Analysts compress hours of repetitive log reading and report drafting into minutes. This acceleration introduces a stark asymmetry: investigative velocity rises while evidentiary verifiability erodes.

Courtroom admissibility under ISO/IEC 27037, ISO/IEC 27042, Daubert, and FRE 702(d) demands byte-level provenance, reproducible execution, and case-specific reliability~\cite{iso27037_2012, iso27042_2015, daubert1993, fre702_2023}. Autoregressive transformers instead compute statistically probable token continuations shaped by prompt phrasing, quantization rounding, and parallel floating-point scheduling. Omitting one exculpatory event due to attention decay or failing to challenge an examiner's unsupported premise~\cite{trivedi2026_when} risks wrongful attribution. Fluency creates a dangerous illusion of forensic certainty.

\subsection{Identified Literature Gaps}
Our systematic audit identifies three structural deficits requiring coordinated community intervention:

\textit{Scarcity of Standardized, Noise-Realistic Forensic Benchmarks:} As audited in Section~\ref{sec:bibliometrics} and Table~\ref{tab:qa_audit}, published evaluations rely heavily on synthetic scenarios, isolated testbeds, or closed institutional casework~\cite{jelodar2026_large, michelet2024_chatgpt, hammar2026_incident}, whereas benchmarking on reference disk corpora such as NIST CFReDS~\cite{nist_cfreds} remains uncommon. Clean high-signal prompts mask real-world failure modes including long-context attention degradation, fragmented database recovery, and elevated false-negative rates under multi-terabyte noise. Initiatives such as ForensicsData~\cite{chakir2025_forensicsdata}, public 13~GB CTF graph benchmarks~\cite{khatiwala2025_evaluating}, and CFTT-inspired timeline evaluation suites~\cite{studiawan2025_towards} represent early progress, yet cross-model validation on realistic forensic images remains incomplete.

\textit{Absence of Certification Protocols for Probabilistic Forensic Tools:} Existing tool validation programs, exemplified by the NIST Computer Forensics Tool Testing (CFTT) program~\cite{nist_cftt}, verify deterministic operations such as write-blocking, disk imaging, and file carving. No harmonized certification baseline yet evaluates non-deterministic generative tools under ISO/IEC 27041~\cite{iso27041_2015}, FRE 702(d)~\cite{fre702_2023}, or the high-risk conformity assessments mandated by Annex~III of the EU AI Act~\cite{eu_ai_act_2024, scanlon2026_investigation, studiawan2025_towards, tiwari2026_xforensicnet}. Equally acute is our QA2 finding: zero of the six audited open-access empirical studies published complete reproducibility manifests.

\textit{Underdeveloped Defenses Against Evidence-Borne Prompt Injection:} Although indirect prompt injection is recognized as a severe forensic hazard~\cite{chernyshev2026_large, kim2026_evidenceborne}, empirical adversarial stress-testing was absent across all six audited open-access empirical studies (0/6) and addressed directly by only one dedicated study in the entire $N=44$ corpus~\cite{kim2026_evidenceborne}. Most academic prototypes concatenate raw artifact text directly into system prompts without pre-inference deterministic filtering or dual-agent auditing.

\subsection{Strategic Roadmap for Future Research}
Transitioning language models from fragile experimental prototypes to court-defensible forensic instruments demands progress along three engineering fronts:

\textit{Hardware-Attested Cryptographic Inference Ledgers:} Next-generation forensic workstations should bind execution manifest $\mathcal{M}$ (Eq.~\ref{eq:forensic_manifest}) to Trusted Platform Modules (TPMs) or Hardware Security Modules (HSMs). Cryptographically signing model checkpoint hashes, decoding parameters, evidence digests, and raw output streams satisfies both ISO/IEC 27042 auditability and EU AI Act Article~12 automatic logging requirements, enabling targeted deterministic CPU replay of disputed courtroom findings.

\textit{Deterministic-First Hybrid Architectures:} Forensic software must enforce a strict architectural boundary in which deterministic parsers (\texttt{plaso}, Volatility, YARA, Sigma, Code Property Graphs) establish immutable timelines, binary offsets, and rule alerts before any air-gapped SLM performs semantic summarization. Pairing canonical sanitization, EBNF grammar constraints, and secondary verifier agents bounds adversarial hijacking without sacrificing triage utility.

\textit{Quantitative Premise-Audit and Impartiality Benchmarks:} Addressing the epistemic failure identified by Trivedi~\cite{trivedi2026_when} requires automated counter-hypothesis prompting protocols that force local models to evaluate exculpatory alternatives and explicitly flag unsupported premises prior to human review~\cite{ko2026_humansupervised}. Establishing standardized CFTT-aligned adversarial test suites containing injected log strings and timestomped artifacts will enable objective cross-laboratory certification.

\subsection{Threats to Validity}
Methodological transparency requires explicitly bounding the validity of this review. Regarding \textit{internal validity and retrieval completeness}, primary programmatic harvesting relied on the Crossref REST API across 28 targeted queries capped at \texttt{rows=35} through November 10, 2026. Although subsequent cross-validation against the open scholarly index OpenAlex achieved a 44/44 (100\%) DOI match and closed Crossref's metadata gaps (raising affiliation visibility to 43/44 [97.7\%] and abstract/full-text coverage to 40/44 [90.9\%]), Crossref's initial relevance-ranked truncation cannot guarantee absolute recall across closed proprietary indexes such as Scopus or Web of Science (for reference, a broad unscreened OpenAlex keyword query for \texttt{"digital forensics" "large language model"} across 2021--2026 returns 2,457 raw records spanning preprints, non-digital domains, and unrefereed grey literature). Regarding \textit{coding reliability}, title screening and eligibility filtering were conducted by a single human author without a second independent human screener; to audit coding consistency transparently, an independent deterministic rule-based second pass over titles and the 40 reconstructed abstracts/full texts yielded 94.1\% agreement ($\kappa=0.821$) on dedicated local-model classification, and all scripts, raw snapshots, and coding matrices are publicly archived on GitHub for third-party verification. Regarding \textit{full-text QA scope}, paywalled full texts were excluded from Table~\ref{tab:qa_audit} rather than extrapolated from abstracts, bounding formal QA1--QA4 scores to the six audited open-access empirical PDFs. Regarding \textit{scope and construct validity}, retaining Axis B studies (Forensics of LLM Systems, $n=5$) alongside two domain-adjacent empirical works (crime-scene image triage~\cite{farber2025_scpaiscp} and microservice incident response~\cite{tian2026_gala}) broadens interdisciplinary insight but requires readers to interpret aggregate counts strictly through the decoupled orthogonal breakdown in Fig.~\ref{fig:subdomain_matrix}. Regarding \textit{language, publication, and temporal validity}, restricting inclusion to English-language archival venues excludes regional non-English forensic scholarship, positive-result publication bias likely underreports silent hallucination failures in operational police work, and vendor deprecation of commercial snapshots (such as early GPT-3.5 and GPT-4 endpoints evaluated in 2023--2024) permanently prevents exact baseline replication. This temporal fragility underscores the central conclusion of this study: forensic science cannot build evidentiary reliability on external, ephemeral commercial infrastructure.
